%% kapa: exam=1 round=25 year=2014 period=1 form=A subject=civil_law subject_ko=민법 seq=17 number=17
%% source: Q-Net 감정평가사 기출문제 (cst00309) artlSeq=5101207 file=extracted/1008/5101207/2014년 제25회 감정평가사 1차시험 기출문제/1교시(A형).hwp
%% conversion: kordoc 4.24.51 -> tools/gen1.py (auto); review=proofread-sonnet; flags=none
\begin{question}{17}
甲은 乙 소유의 X토지를 매수하면서 자신과 명의신탁약정을 한 친구 丙 앞으로 소유권이전등기를 해 주기로 乙과 합의하였고, 그에 따라 X토지의 소유권이전등기가 丙 명의로 마쳐졌다. 다음 설명 중 옳지 않은 것은? (다툼이 있으면 판례에 의함)
\begin{choices}
  \item 丙 명의의 소유권이전등기는 무효이다.
  \item 乙은 丙에게 진정명의회복을 원인으로 하는 소유권이전등기를 청구할 수 있다.
  \item 甲은 乙을 대위하여 丙에게 소유권이전등기의 말소청구를 할 수 없다.
  \item 乙은 甲에 대하여 소유권이전등기의무를 부담한다.
  \item 丙으로부터 정당하게 X토지를 매수한 丁이 甲과 丙 사이의 명의신탁 약정사실을 알았더라도, 丁은 丙으로부터 X토지를 유효하게 취득할 수 있다.
\end{choices}
\end{question}
\answer{3}
